\documentclass[10pt,a4paper]{amsart}
\usepackage[margin=19mm]{geometry}
\usepackage{amsmath,amssymb,mathtools}
\usepackage[scr=rsfs,cal=euler]{mathalfa}
\usepackage{xfrac}
\usepackage[unicode,hidelinks,bookmarks=false]{hyperref}
\newcommand{\ie}{i.e.\ }
\newcommand{\cf}{cf.\ }
\newcommand{\resp}{respectively\ }
\newcommand{\Subsection}{\subsection}
\providecommand{\nobreakdash}{\nobreak\hbox{-}}
\setlength{\emergencystretch}{2em}
\multlinegap0pt
\usepackage[shortlabels]{enumitem}
\usepackage{bbm}
\usepackage{mathtools}
\usepackage[normalem]{ulem}
\usepackage{tikz}
\usepackage{cancel}
\usepackage{siunitx}
\usepackage{multirow}
\usepackage{amssymb}
\usepackage{booktabs}
\newtheorem{prop}{Proposition}
\newtheorem{lemma}{Lemma}
\DeclareMathOperator*{\argmax}{argmax}
\DeclareMathOperator*{\argmaxlex}{\overset{\lex}{\argmax}}
\DeclareMathOperator*{\argmaxxlex}{\overset{[2],\lex}{\argmax}}
\DeclareMathOperator{\lex}{lex}
\DeclareMathOperator*{\maxlex}{\overset{\lex}{\max}}
\DeclareMathOperator{\perfV}{perf_V}
\title{Optimal matching under size priority}
\author{Nathana\"el Enriquez, Mike Liu, Laurent M\'enard, Vianney Perchet}
\date{Source: arXiv:2601.20502v1 (CC BY 4.0)}
\begin{document}
\maketitle

\noindent\emph{Source and licence.} Optimal matching under size priority. Version arXiv:2601.20502v1 (CC BY 4.0), distributed by arXiv under the Creative Commons Attribution 4.0 International licence, http://creativecommons.org/licenses/by/4.0/, as recorded in the arXiv licence metadata for this exact version and shown on the abstract page https://arxiv.org/abs/2601.20502v1. Full licence notice: \texttt{licenses/arxiv-2601.20502-CC-BY-4.0.txt}.\par\smallskip

\noindent\textit{Editorial scope.} Counted unit: the article's prop prop:uniqueness (source0.tex lines 1742--1745). The article's complete original proof of it is delivered with it (source0.tex lines 1747--1898, one \texttt{proof} environment of 152 lines). No mathematics is written, repaired or re-derived: the statement and the proof are the article's own lines, in the article's own notation, and no retained line is substituted or re-flowed.  Retained uncounted as the source-local context the proof needs: lines 1138--1147; lines 1336--1362; lines 1569--1571; lines 1610--1648.  One declared adaptation: exactly 2 ancillary strings inside the retained spans are omitted (image-inclusion commands and, where the source repeats a label, the repeated label definitions) and no other line differs from the source; the figure environments, with their placement, captions and labels, are kept, and the package records the omitted strings in fidelity receipts bound to the affected bodies.  No retained line contains a citation, so the wrapper carries no bibliography and no bibliography entry.  The retained lines contain the article's own diagram code, which the wrapper typesets with the standard tikz packages; no image file is delivered and the package declares no asset dependency.  Editorial formatting only, in addition: an AMS \texttt{amsart} preamble that restates the article's own declarations and macros used by the retained lines, namely the environments \texttt{prop}, \texttt{lemma} on separate counters, together with the standard packages those lines need.  The retained environments print in this document as Proposition 1, Lemma 1 in order of appearance, whereas the article prints the same items with the numbering of its own full document.  The complete native source of the article as distributed in the arXiv e-print for this exact version is bundled unchanged as \texttt{sources/193525/source0.tex}.
\begin{equation}\label{eq:systemzoomh0singular}
\left\{
\begin{aligned}
    \forall 0 < j \leq k, h_j(t)&= \hat{\phi}(1-\mathbb{E}[h_{{k-j}}(W-t)]), \\
    h_0(t)&= \mathbbm{1}_{t \geq 0} \hat{\phi}(1-\mathbb{E}[h_{k}(W-t)]),\\
    \forall 0 \leq j \leq k, \lim_{t \rightarrow -\infty} h_j(t) &< \lim_{t \rightarrow + \infty}h_j(t),  \\
    \forall 0 \leq j < k, \lim_{t \rightarrow + \infty}h_j(t) &= \lim_{t \rightarrow -\infty}h_{j+1}(t), 
\end{aligned}
\right.
\end{equation}

\begin{prop}\label{prop:Zconstruction}
    Let $\mathbb{T}$ be a unimodular decorated BGW tree with reproduction law $\pi$ and edge weight law $\omega$. Let $\mathcal{H}$ be a solution to the system given by System~\eqref{eq:systemzoomh0singular} with input $\pi$ and $\omega$.
    \begin{enumerate}[label=\roman*)]
        \item There exists a random decorated tree $(\mathbb{T}',((i_{\mathcal{H}},Z_{\mathcal{H}})(u,v))_{(u,v) \in \overset{\rightarrow}{E}})$ such that the marginal law of each $(i_{\mathcal{H}},Z_{\mathcal{H}})(u,v)$ is $\zeta'_{\mathcal{H}}$, the marginal law of $\mathbb{T}'$ is the law of $\mathbb{T}$ and such that for any $(u,v) \in \overset{\rightarrow}{E}$:
        \begin{equation}\label{eq:systemrecursionzoombis}
    \begin{aligned}
        (i_{\mathcal{H}},Z_{\mathcal{H}})(u,v)=\maxlex\left( (0,0) , \maxlex_{\substack{ u' \sim v \\ u' \neq u}}\left( (k,w(v,u'))-(i_{\mathcal{H}},Z_{\mathcal{H}})(v,u') \right) \right).
    \end{aligned}
\end{equation}
        \item The rule
        \begin{align}\label{eq:decisionrule}
            (u,v) \in \mathbb{M}_{\mathcal{H}} \Leftrightarrow (i_{\mathcal{H}},Z_{\mathcal{H}})(u,v)+(i_{\mathcal{H}},Z_{\mathcal{H}})(v,u)\overset{\lex}{<} (k,w(u,v)) 
        \end{align}
        defines a matching on $\mathbb{T}$. This rule is equivalent to the following vertex rule:
        $u$ is matched if and only if there exists $v\sim u$ such that
\begin{equation}\label{eq:vertexrulebis}
        \vphantom{\underset{v' \sim u}{\overset{\lex}{\argmax}}}
        (k,w(u,v))-(i_{\mathcal{H}},Z_{\mathcal{H}})(u,v) \overset{\lex}{>} (0,0),
\end{equation}
in which case $u$ is matched to the vertex $v$ defined by
\begin{equation}
        v=\underset{v' \sim u}{\overset{\lex}{\argmax}} \left( (k,w(u,v'))-(i_{\mathcal{H}},Z_{\mathcal{H}})(u,v') \right).
\end{equation}

        \item The random decorated tree $(\mathbb{T}', (i_{\mathcal{H}},Z_{\mathcal{H}})_{(u,v) \in \overset{\rightarrow}{E}})$ viewed as $\mathbb{T}'$ with an additional decoration is unimodular.
        \end{enumerate}
\end{prop}

\begin{prop}\label{prop:optimality}
    The matching  $\mathbb{M}_{\mathcal{H}}$ defined in Proposition~\ref{prop:Zconstruction} is optimal on $\mathbb{T}$. Furthermore, for any coupling $(\mathbb{T},\mathbb{M}_{\mathcal{H}},\mathbb{M})$ with $\mathbb{M}$ an optimal matching on $\mathbb{T}$, $\mathbb{M}$ and $\mathbb{M}_{\mathcal{H}}$ share the same set of unmatched vertices.
\end{prop}

By hypothesis, $v_{-1} \neq v_{1}$ hence by applying $n$, $n(v_{-1})=v_{-2} \neq v_0=n(v_1)$, we deduce that 
\begin{align*}
    (i^s_{\mathcal{H}},Z^s_{\mathcal{H}})(v_{-2},v_{-1})=(j_{\mathcal{H}}^s,w_{\mathcal{H}}^s)(v_{-1},v_{0})-(i^s_{\mathcal{H}},Z^s_{\mathcal{H}})(v_{-1},v_0).
\end{align*}
Now let us look at $(i^s_{\mathcal{H}},Z^s_{\mathcal{H}})(v_{-1},v_0)$:
\begin{equation}
\begin{aligned}\label{eq:aldous2}
    (i^s_{\mathcal{H}},Z^s_{\mathcal{H}})(v_{-1},v_0) &= \underset{{ y \sim v_0}}{\maxlex}\left( (j_{\mathcal{H}}^s,w_{\mathcal{H}}^s)(v_0,y)-(i^s_{\mathcal{H}},Z^s_{\mathcal{H}})(v_0,y) \right). \\
    &\geq (j_{\mathcal{H}}^s,w_{\mathcal{H}}^s)(v_0,v_1)-(i^s_{\mathcal{H}},Z^s_{\mathcal{H}})(v_0,v_1).
\end{aligned}
\end{equation}
Putting everything together:
\begin{equation}\label{eq:aldous4}
\begin{aligned}
    &\left((j_{\mathcal{H}}^s,w_{\mathcal{H}})(v_{-1},v_0)-(j_{\mathcal{H}}^s,w_{\mathcal{H}}^s)(v_{1},v_0) \right)\mathbbm{1}_{v_{-1} \neq v_1} \\
    =&\left(((i^s_{\mathcal{H}},Z^s_{\mathcal{H}})(v_{-2},v_{-1})+(i^s_{\mathcal{H}},Z^s_{\mathcal{H}})(v_{-1},v_0)-(j_{\mathcal{H}}^s,w_{\mathcal{H}}^s)(v_{1},v_0)\right)\mathbbm{1}_{v_{-1} \neq v_1} \\
    =&\bigg((i^s_{\mathcal{H}},Z^s_{\mathcal{H}})(v_{-2},v_{-1})+(i^s_{\mathcal{H}},Z^s_{\mathcal{H}})(v_0,v_1) \\
    +&\bigg[(i^s_{\mathcal{H}},Z^s_{\mathcal{H}})(v_{-1},v_0)-(i^s_{\mathcal{H}},Z^s_{\mathcal{H}})(v_0,v_1)-(j_{\mathcal{H}}^s,w_{\mathcal{H}}^s)(v_{1},v_0)\bigg]\bigg)\mathbbm{1}_{v_{-1} \neq v_1} \\
    \overset{\lex}{\geq}&\left((i^s_{\mathcal{H}},Z^s_{\mathcal{H}})(v_{-2},v_{-1}) -(i^s_{\mathcal{H}},Z^s_{\mathcal{H}})(v_{0},v_{1})\right)\mathbbm{1}_{v_{-1} \neq v_1} .
\end{aligned}
\end{equation}
Now decompose the weights depending on whether $v_{-1}=v_0$ or $v_{1}=v_0$ and use the fact that $j_{\mathcal{H}}^s(u,v)=k$  if $u \neq v$:
\begin{align*}
    &\left((j_{\mathcal{H}}^s,w_{\mathcal{H}})(v_{-1},v_0)-(j_{\mathcal{H}}^s,w_{\mathcal{H}}^s)(v_{1},v_0) \right)\mathbbm{1}_{v_{-1} \neq v_1} \\
    =&(k,w(v_{-1},v_{0}))\mathbbm{1}_{v_{-1}\neq v_{0}}-(k,w(v_1,v_{0}))\mathbbm{1}_{v_1 \neq v_{0}} \\
    +&(j_{\mathcal{H}}^s,w_{\mathcal{H}}^s)(o,o)\mathbbm{1}_{v_{-1}=v_{0}}-(j_{\mathcal{H}}^s,w_{\mathcal{H}}^s)(o,o)\mathbbm{1}_{v_{1}=v_{0}}.
\end{align*}
Rearranging the different terms, and recalling that we are working on the event $v_1 \neq v_{-1}$, we obtain:
\begin{align*}
    &\left((\mathbbm{1}_{v_{-1}\neq v_0}-\mathbbm{1}_{v_1\neq v_0}),\left(w^s(v_{-1},v_0)\mathbbm{1}_{v_{-1}\neq v_0}-w^s(v_{1},v_0)\mathbbm{1}_{v_1 \neq v_0}\right) \right)\mathbbm{1}_{v_{-1}\neq v_1} \\
    &+ \left((j_{\mathcal{H}}^s, w_{\mathcal{H}}^s)(o,o)  \mathbbm{1}_{v_0=v_{-1}, v_0 \neq v_1}-( j_{\mathcal{H}}^s, w_{\mathcal{H}}^s)(o,o)  \mathbbm{1}_{v_0=v_{1}, v_0 \neq v_{-1}} \right) \\
    &\overset{\lex}{\geq} \left( (i^s_{\mathcal{H}},Z^s_{\mathcal{H}})(v_{-2},v_{-1}) -(i^s_{\mathcal{H}},Z^s_{\mathcal{H}})(v_{0},v_{1})\right)\mathbbm{1}_{v_{-1} \neq v_1}.
\end{align*}
So we set
\begin{align*}
    X'&:=\left[ (j_{\mathcal{H}}^s, w_{\mathcal{H}}^s)(o,o)  \mathbbm{1}_{o=n_{\mathcal{H}}(o),o\neq n(o)}- (j_{\mathcal{H}}^s, w_{\mathcal{H}}^s)(o,o)  \mathbbm{1}_{o=n_{\mathcal{H}}(o),o\neq n(o)}\right], \\
    Y&:=(i^s_{\mathcal{H}},Z^s_{\mathcal{H}})(v_{-2},v_{-1})\mathbbm{1}_{v_{-1} \neq v_1} , \\
    Y'&:=(i^s_{\mathcal{H}},Z^s_{\mathcal{H}})(v_{0},v_{1})\mathbbm{1}_{v_{-1} \neq v_1} .
\end{align*}

\begin{prop}\label{prop:uniqueness}
    Let $\mathbb{M}$ be an optimal matching on $\mathbb{T}$ and let $(\mathbb{T},\mathbb{M},\mathbb{M}_{\mathcal{H}})$ be a coupling where $\mathbb{M}_{\mathcal{H}}$ is as in Proposition~\ref{prop:Zconstruction}.
    Then $\mathbb{M}=\mathbb{M}_{\mathcal{H}}$ almost surely.
\end{prop}

\begin{proof}
We will show that almost surely, $\mathbb{M}^s=\mathbb{M}_{\mathcal{H}}^s$ and deduce that almost surely, $\mathbb{M}=\mathbb{M}_{\mathcal{H}}$ through the projection that forgets self-loops. In the whole proof, let $A$ be the event when $n_{\mathcal{H}}(o)\neq n(o)$ and assume $\mathbb P(A) \neq 0$.


Since $\mathbb{M}$ is optimal, we know from the previous proof that $\mathbb{M}$ and $\mathbb{M}_{\mathcal{H}}$ share the same set of unmatched vertices. 
Recalling Equation~\eqref{eq:aldous4}, we get
\begin{align*}
&\left((j_{\mathcal{H}}^s,w_{\mathcal{H}}^s)(v_{-1},v_0)-(j_{\mathcal{H}}^s,w_{\mathcal{H}}^s)(v_{1},v_0) \right)\mathbbm{1}_{v_{-1} \neq v_1} \\
-&\left((i^s_{\mathcal{H}},Z^s_{\mathcal{H}})(v_{-2},v_{-1}) -(i^s_{\mathcal{H}},Z^s_{\mathcal{H}})(v_{0},v_{1})\right)\mathbbm{1}_{v_{-1} \neq v_1} \\
\overset{\lex}{\geq}&\left((j_{\mathcal{H}}^s,w_{\mathcal{H}}^s)(v_0,v_1)-(i^s_{\mathcal{H}},Z^s_{\mathcal{H}})(v_0,v_1)-(i^s_{\mathcal{H}},Z^s_{\mathcal{H}})(v_{-1},v_0)\right)\mathbbm{1}_{v_1 \neq v_{-1}} \\
\overset{\lex}{\geq} &\,0.
\end{align*}
The previous subsection also shows that
\[
\mathbb{E}[\left((j_{\mathcal{H}}^s,w_{\mathcal{H}}^s)(v_0,v_1)-(i^s_{\mathcal{H}},Z^s_{\mathcal{H}})(v_0,v_1)-(i^s_{\mathcal{H}},Z^s_{\mathcal{H}})(v_{-1},v_0)\right)\mathbbm{1}_{v_1 \neq v_{-1}}]
\]
exists and is equal to zero. Indeed, the expectation of the difference of performances is 
\[
\perfV(\mathbb{M}_{\mathcal{H}})-\perfV(\mathbb{M})=(0,0),
\]
and we deduce that $\left((j_{\mathcal{H}}^s,w_{\mathcal{H}}^s)(v_0,v_1)-(i^s_{\mathcal{H}},Z^s_{\mathcal{H}})(v_0,v_1)-(i^s_{\mathcal{H}},Z^s_{\mathcal{H}})(v_{-1},v_0)\right)\mathbbm{1}_{v_1 \neq v_{-1}}$ is integrable and that
\begin{align*}
\mathbb{E}\left[((j_{\mathcal{H}}^s,w_{\mathcal{H}}^s)(v_0,v_1)-(i^s_{\mathcal{H}},Z^s_{\mathcal{H}})(v_0,v_1)-(i^s_{\mathcal{H}},Z^s_{\mathcal{H}})(v_{-1},v_0))\mathbbm{1}_{v_1 \neq v_{-1}} \right]=(0,0) .
\end{align*}
Since Equation~\ref{eq:aldous2} implies that this variable is almost surely bigger than $(0,0)$, we deduce that it is equal to $(0,0)$ almost surely, hence that
\begin{align}\label{eq:Aldous5}
(i^s_{\mathcal{H}},Z^s_{\mathcal{H}})(v_{-1},o) =(j_{\mathcal{H}}^s,w_{\mathcal{H}}^s)(o,v_1)-(i^s_{\mathcal{H}},Z^s_{\mathcal{H}})(o,v_1).
\end{align}
By definition of $n^s_{\mathcal{H}}$,
\begin{align}\label{eq:Aldous6}
v_{-1}=n^s_{\mathcal{H}}(v_0) = \overset{\lex}{\argmax}_{y \sim v_0 }((j_{\mathcal{H}}^s,w_{\mathcal{H}}^s)(o,y)-(i^s_{\mathcal{H}},Z^s_{\mathcal{H}})(o,y)).
\end{align}
By the recursive equation of the messages: 
\begin{align}\label{eq:Aldous7}
    (i^s_{\mathcal{H}},Z^s_{\mathcal{H}})(v_{-1},o)=\maxlex_{\substack {y \sim o  y \neq v_{-1}}}((j_{\mathcal{H}}^s,w_{\mathcal{H}}^s)(o,y)-(i^s_{\mathcal{H}},Z^s_{\mathcal{H}})(o,y))    
\end{align}
Combining Equations~\eqref{eq:Aldous5} and \eqref{eq:Aldous7} then using Equation~\eqref{eq:Aldous6}, we get that $v_1$ achieves the maximum among the list of $(j_{\mathcal{H}}^s,w_{\mathcal{H}}^s)(o,y)-(i^s_{\mathcal{H}},Z^s_{\mathcal{H}})(o,y)$ stripped of its maximum. 
In other words, $v_1=n(o)$ achieves the second maximum of $(j_{\mathcal{H}}^s,w_{\mathcal{H}}^s)(o,y)-(i^s_{\mathcal{H}},Z^s_{\mathcal{H}})(o,y)$ when $y$ spans the neighbours of $o$. We will denote this quantity by ${\underset{y \sim o}{\argmaxxlex}}( (j_{\mathcal{H}}^s,w_{\mathcal{H}}^s)(o,y)-(i^s_{\mathcal{H}},Z^s_{\mathcal{H}})(o,y))$. 
Thus:
\[\mathbb{P}\left( n(o)=\underset{y \sim o}{\argmaxlex}((j_{\mathcal{H}}^s,w_{\mathcal{H}}^s)(o,y)-(i^s_{\mathcal{H}},Z^s_{\mathcal{H}})(o,y)) \text{ or } \underset{y \sim o}{\argmaxxlex}((j_{\mathcal{H}}^s,w_{\mathcal{H}}^s)(o,y)-(i^s_{\mathcal{H}},Z^s_{\mathcal{H}})(o,y))  \right)=1. \]
By unimodularity, this must be true for any vertex $u \in V(\mathbb{T})$. In particular, we obtain that
\[ n(u)\neq n_{\mathcal{H}}(u)  \implies n(u)= \bigg(\argmaxxlex_{v \sim u} (j_{\mathcal{H}}^s,w_{\mathcal{H}}^s)(u,v)-(i_{\mathcal{H}}^s,Z_{\mathcal{H}}^s)(u,v) \bigg) . \]
In addition, when an edge of the matching is selected according to the second largest value, this rule has to be involutive, meaning:
\begin{align} \label{eq:argmax2involutif}
    n(u)\neq n_{\mathcal{H}}(u)  \implies &\left\{  n(u)= \bigg(\argmaxxlex_{v \sim u} (j_{\mathcal{H}}^s,w_{\mathcal{H}}^s)(u,v)-(i_{\mathcal{H}}^s,Z_{\mathcal{H}}^s)(u,v) \bigg) \right. \notag \\
    & \qquad  \left. \text{ and } u= \bigg(\argmaxxlex_{v \sim n(u)} (j_{\mathcal{H}}^s,w_{\mathcal{H}}^s)(n(u),v)-(i_{\mathcal{H}}^s,Z_{\mathcal{H}}^s)(n(u),v) \bigg) \right\}.
\end{align}
The remainder of the argument will be based on the fact that it is impossible to use the second largest rule to define a matching in a unimodular fashion.

Let us set
\begin{align*}
\tilde{v}_0&:=o ,\\
\tilde{v}_{1}&:=\underset{y \sim o}{\argmaxxlex}((j_{\mathcal{H}}^s,w_{\mathcal{H}}^s)(o,y)-(i^s_{\mathcal{H}},Z^s_{\mathcal{H}})(o,y)) &&\text{if }o\text{ is not isolated} ,\\
\tilde{v}_{2}&:=\underset{y \sim \tilde{v}_1}{\argmaxlex}((j_{\mathcal{H}}^s,w_{\mathcal{H}}^s)(\tilde{v}_1,y)-(i^s_{\mathcal{H}},Z^s_{\mathcal{H}})(\tilde{v}_1,y)) &&\text{if }o\text{ is not isolated} ,\\
\tilde{v}_{3}&:=\underset{y \sim \tilde{v}_2}{\argmaxxlex}((j_{\mathcal{H}}^s,w_{\mathcal{H}}^s)(\tilde{v}_2,y)-(i^s_{\mathcal{H}},Z^s_{\mathcal{H}})(\tilde{v}_2,y))&&\text{if }o\text{ is not isolated} ,\\
\tilde{v}_{-1}&:=\underset{y \sim o}{\argmaxlex}((j_{\mathcal{H}}^s,w_{\mathcal{H}}^s)(o,y)-(i^s_{\mathcal{H}},Z^s_{\mathcal{H}})(o,y)) ,\\
\tilde{v}_{-2p}&:=\underset{y \sim w_{-2p+1}}{\argmaxxlex}((j_{\mathcal{H}}^s,w_{\mathcal{H}}^s)(\tilde{v}_{-2p+1},y)-(i^s_{\mathcal{H}},Z^s_{\mathcal{H}})(\tilde{v}_{-2p+1},y))  &&\forall p \in \mathbb{N}^*,\\
\tilde{v}_{-2p-1}&:=\underset{y \sim w_{-2p}}{\argmaxlex}((j_{\mathcal{H}}^s,w_{\mathcal{H}}^s)(\tilde{v}_{-2p},y)-(i^s_{\mathcal{H}},Z^s_{\mathcal{H}})(\tilde{v}_{-2p},y))  &&\forall k \in \mathbb{N}^* .
\end{align*}
See Figure~\ref{fig:infpath} for an illustration.
The advantage of the $\tilde{v}_p$ instead of the previously used $v_p$ is that they no longer depend on $\mathbb{M}$ which we do not have much control over. Instead, they depend on the messages directly.
\begin{figure}
    \centering
    
    \caption{Definition of the path $(\tilde{v}_p)_{p\leq3}$.}
    \label{fig:infpath}
\end{figure}


We define the following events  for $p<0$ and $N < 0$
\begin{align}
A_{2p}^{(1)} &:= \left\{ \tilde{v}_{2p} \neq \tilde{v}_{2p+1} \neq \tilde{v}_{2p-1} \right\},\\
A_{2p}^{(2)} &:= \left\{  \tilde{v}_{2p+1} =\argmaxxlex_{u \sim \tilde{v}_{2p}}\bigg( (j_{\mathcal{H}}^s,w_{\mathcal{H}}^s)(\tilde{v}_{2p},u)- (i_{\mathcal{H}}^s,Z_{\mathcal{H}}^s)(\tilde{v}_{2p},u)\bigg) \right\},\\
A_{2p} &:= A_{2p}^{(1)} \cap A_{2p}^{(2)},\\
C_{N} &:=\bigcap_{N \leq 2p < 0} A_{2p} . 
\end{align}
The first event $A_{2p}^{(1)}$ corresponds to the fact that the portion of the path around $\tilde{v}_{2p}$ is injective and  $A_{2p}^{(2)}$ corresponds to \eqref{eq:argmax2involutif}. See figure~\ref{fig:uniquenessevent}.

\begin{figure}
    \centering
    
    \caption{Definition of the event $A_{2p}$.}
    \label{fig:uniquenessevent}
\end{figure}

By construction, we have $(\tilde{v}_{-2p},\tilde{v}_{-2p-1})\in \mathbb{M}^s_{\mathcal{H}}$ for any $p\geq0$. Recall that we are interested in the event
\begin{equation}
    A := \{n_{\mathcal{H}}(o)\neq n(o)\}.
\end{equation}
Note that on the event $A$ we have
\[
\tilde{v}_{-2} = v_{-2}, \, \tilde{v}_{-1} = v_{-1}, \,\tilde{v}_{0} = v_0 \text{ and } \tilde{v}_1 = v_1.
\]
In addition we have
\[\mathbb{P}\left(\bigcup_{p \geq 0}\left\{ (\tilde{v}_{-2p-1},\tilde{v}_{-2p-2})\in \mathbb{M}\right\} \middle| A \right)=1. \]

Furthermore, by the second part of Proposition~\ref{prop:optimality}, almost surely, for any $v\in V$, $n_{\mathcal{H}}(v)=v$ if and only if $n(v)=v$. From this, we deduce that on $A$ we have $\tilde{v}_{-1}\neq \tilde{v}_0$ and, by induction, $  \forall k \in \mathbb{N}, \tilde{v}_{-k} \neq \tilde{v}_{-k-1}$ (no loops in the path) and also that $\forall k \in \mathbb{N}, \tilde{v}_{-k-2} \neq \tilde{v}_{-k}$ (the path cannot go back up). This holds similarly for $\tilde{v}_1$,$\tilde{v}_2$ and $\tilde{v}_3$. In short, we showed that $A$ is contained in the event where the path $(\tilde{v}_k)_{k\leq 3}$ is infinite and injective.



For the root and $\tilde{v}_2$, a direction reversal occurs, recall
\[ \tilde{v}_1=\argmaxxlex_{u \sim o} \bigg( (j_{\mathcal{H}}^s,w_{\mathcal{H}}^s)(o,u)- (i_{\mathcal{H}}^s,Z_{\mathcal{H}}^s)(o,u)\bigg)
.  \]
We will then set
\begin{align}
\tilde{A}_0 &:= \left\{ \tilde{v}_0=\argmaxxlex_{u \sim \tilde{v}_1} \bigg( (j_{\mathcal{H}}^s,w_{\mathcal{H}}^s)(\tilde{v}_1,u)- (i_{\mathcal{H}}^s,Z_{\mathcal{H}}^s)(\tilde{v}_1,u)\bigg), \tilde{v}_0\neq \tilde{v}_{-1}, \tilde{v}_ \neq \tilde{v}_1 
\right\} , \\
\tilde{A}_2 &:=\left\{ \tilde{v}_2=\argmaxxlex_{u \sim \tilde{v}_3} \bigg( (j_{\mathcal{H}}^s,w_{\mathcal{H}}^s)(\tilde{v}_3,u)- (i_{\mathcal{H}}^s,Z_{\mathcal{H}}^s)(\tilde{v}_3,u)\bigg) , \tilde{v}_2\neq \tilde{v}_1, \tilde{v}_3 \neq \tilde{v}_2 \right\} .
\end{align}

For every $N <0$ we have
\[ \mathbb{P}(A) \leq \mathbb{P}(\tilde A_0\cap \tilde A_2 \cap C_N) = \mathbb{P}\left( \tilde A_0 \cap \tilde A_2 \cap \bigcap_{N\leq 2p < 0} A_{2p} \right)  . \]
To conclude the argument, we want to show that the right hand side has probability zero when $N \rightarrow -\infty$. We will use the following Lemma, whose proof is postponed at the end of the section.
\begin{lemma}\label{lem:doublyinfinitepath}
Fix some $H>2$, then $\mathbb{P}(C_{-2H-2}\cap \tilde{A}_0)=\mathbb{P}(C_{-2H} \cap \tilde{A}_0 \cap \tilde{A}_2)$.
\end{lemma}

Now, if we assume that $ \lim_{N \rightarrow -\infty} \mathbb{P}(C_{N}\cap \tilde{A}_0 \cap \tilde{A_2})>0$, taking $H \rightarrow \infty$ in Lemma \ref{lem:doublyinfinitepath} yields
\begin{equation}\label{eq:infinitepath}
 \mathbb{P}\left(  \tilde{A}_{2}  \bigg| \tilde A_0 \cap \bigcap_{p' < 0} A_{2p'}  \right) =1.
\end{equation}

Assume that $\tilde{v}_2\neq \tilde{v}_3$. The only dependency between the messages and variables in the connected component $\mathbb{T}_{\tilde{v}_{3}}$ of $\tilde{v}_{3} \in \mathbb{T}\setminus \{\tilde{v}_{2} \}$ on the one hand, and its complement $\mathbb T_{\tilde{v}_{2}}$ on the other hand, is through the triplet
\[(i_{\mathcal{H}},Z_{\mathcal{H}})(\tilde{v}_{3},\tilde{v}_{2}),(i_{\mathcal{H}},Z_{\mathcal{H}})(\tilde{v}_{2},\tilde{v}_{3}) ,w(\tilde{v}_{3},\tilde{v}_{2}). \] 
From this we deduce that conditionally on the triplet
\[(i_{\mathcal{H}},Z_{\mathcal{H}})(\tilde{v}_{3},\tilde{v}_{2}),(i_{\mathcal{H}},Z_{\mathcal{H}})(\tilde{v}_{2},\tilde{v}_{3}) ,w(\tilde{v}_{3},\tilde{v}_{2}), \]
 the two decorated subtrees \[\left(\mathbb{T}_{\tilde{v}_{2}},\tilde{v}_{2},(i_{\mathcal{H}},Z_{\mathcal{H}})(u,v)_{(u,v)\in \overset{\rightarrow}{E}(\mathbb{T}_{\tilde{v}_{2}})}\right),\left(\mathbb{T}_{\tilde{v}_{3}},\tilde{v}_{3},(i_{\mathcal{H}},Z_{\mathcal{H}})(u,v)_{(u,v)\in \overset{\rightarrow}{E}(\mathbb{T}_{\tilde{v}_{3}})}\right) \] are independent. Furthermore, according to Equation~\eqref{eq:infinitepath}, conditionally on $\tilde{A_0}\cap \bigcap_{p'<0}A_{2p'}$, $\tilde{A}_{2}$ happens almost surely, which contains the event that $\tilde{v}_2 \neq \tilde{v}_3$.
 
As a consequence, we have that the law of
$\left(\mathbb{T}_{\tilde{v}_{3}},\tilde{v}_{3},(i_{\mathcal{H}},Z_{\mathcal{H}})(u,v)_{(u,v)\in \overset{\rightarrow}{E}(\mathbb{T}_{\tilde{v}_{3}})}\right)$
conditionally on $\tilde A_0 \cap \bigcap_{p' < 0}A_{2p'}$, which is measurable with respect to the other tree, can be condensed through the conditional distribution $\mathcal{P}$ of the triplet  \[(i_{\mathcal{H}},Z_{\mathcal{H}})(\tilde{v}_{3},\tilde{v}_{2}),(i_{\mathcal{H}},Z_{\mathcal{H}})(\tilde{v}_{2},\tilde{v}_{3}) ,w(\tilde{v}_{3},\tilde{v}_{2}) \] under $\tilde A_0 \cap \bigcap_{p' < 0}A_{2p'}$.



We will then show that no matter the values of this triplet, the event $\tilde A_{2}$ will not happen with probability $1$. We thus write that
\begin{align*}
    &\mathbb{P}\left(\overline{\tilde{A_{2}}} \bigg| \tilde A_0 \cap \bigcap_{p' < 0} A_{2p'}\right) \\
    &=\iiint \mathbb{P}\left(\overline{\tilde{A}_{2}} \bigg| \bigcap_{p' <0} A_{2p'}\cap \tilde{A}_0, (i_\mathcal{H},Z_{\mathcal{H}})(\tilde{v}_{2},\tilde{v}_{3})=(a,b), (i_\mathcal{H},Z_{\mathcal{H}})(\tilde{v}_{3},\tilde{v}_{2})=(a',b'), w(\tilde{v}_{2},\tilde{v}_{3})=\alpha  \right)\\
    &\qquad \qquad \qquad \mathrm{d}\mathcal{P}((a,b),(a',b'),\alpha) \\
    &=\iiint \mathbb{P}\left(\overline{\tilde{A}_{2}} \bigg| \tilde{v}_2 \neq \tilde{v}_3, (i_\mathcal{H},Z_{\mathcal{H}})(\tilde{v}_{2},\tilde{v}_{3})=(a,b), (i_\mathcal{H},Z_{\mathcal{H}})(\tilde{v}_{3},\tilde{v}_{2})=(a',b'), w(\tilde{v}_{2},\tilde{v}_{3})=\alpha  \right)\\
    &\qquad \qquad \qquad \mathrm{d}\mathcal{P}((a,b),(a',b'),\alpha )
\end{align*}
where $\mathcal{P}$ is the conditional law of the triplet $(i_{\mathcal{H}},Z_{\mathcal{H}})(\tilde{v}_{3},\tilde{v}_{2}),(i_{\mathcal{H}},Z_{\mathcal{H}})(\tilde{v}_{2},\tilde{v}_{3}) ,w(\tilde{v}_{3},\tilde{v}_{2})$ under $\tilde{A}_0\cap\bigcap_{p' < 0} A_{2p'}$. By definition, we must have that the support of $\mathcal{P}$ is included in the space $(a,b)+(a',b') \overset{\lex}{\geq }(k,\alpha)$ since $(\tilde{v}_{3},\tilde{v}_{2}) \notin \mathbb{M}_h^s$. As equality almost surely does not occur, we almost surely have $(a,b)+(a',b') \overset{\lex}{> }(k,\alpha)$ . The next lemma concludes the proof by showing that being in this space ensures that this conditional density is always positive.
\begin{lemma}\label{lem:Mhstop}
    For every $(a,b),(a',b') \in \mathrm{supp}(\zeta_{h}')$, $ \alpha \in \mathrm{supp} (\omega)$ such that $(a,b)+(a',b')\overset{\lex}{>}(k,\alpha)$, \[ \mathbb{P} \left( \overline{\tilde A_{2}} \middle| \tilde{v}_2 \neq \tilde{v}_3 , \,(i_\mathcal{H},Z_{\mathcal{H}})(\tilde{v}_{2},\tilde{v}_{3})=(a,b), (i_\mathcal{H},Z_{\mathcal{H}})(\tilde{v}_{3},\tilde{v}_{2})=(a',b'), w(\tilde{v}_{2},\tilde{v}_{3})=\alpha \right)    >0   . \]
\end{lemma}
As a conclusion, this shows a contradiction with the assumption that  \[  \mathbb{P}\left(  \bigcap_{p'<0}A_{2p'} \cap \tilde{A_0}\cap \tilde{A_2} \right)>0.\]
We deduce that this probability has to be zero, but since we showed that this probability is bigger than $\mathbb{P}(A)$, we get that $\mathbb{P}(A)=0$, which translates to
\[ \mathbb{P}(n(o)=n_{\mathcal{H}}(o))=1.  \]
By unimodularity, we deduce that almost surely, for all $v \in V(\mathbb{T})$, $n(v)=n_{\mathcal{H}}(v)$, hence $\mathbb{M}=\mathbb{M}_{\mathcal{H}}$.

\end{proof}

\end{document}
